\section{Introduction}\label{sec:intro}

The proximal point algorithm seeks a zero of a set-valued operator by solving an implicit inclusion at each iteration. Martinet's regularization of variational inequalities~\cite{Martinet1970} and Rockafellar's analysis for maximal monotone operators~\cite{Rockafellar1976} established the classical theory. Monotonicity supplies a distance inequality, while maximality supplies the existence of proximal points. Beyond this setting, these two roles must be examined separately: an estimate for a proximal step does not by itself show that such a step exists.

Luke and Tam~\cite{LukeTam2025} study this distinction locally. Their submonotonicity condition quantifies a weakening of monotonicity on a localized operator graph. Together with metric subregularity, local resolvent coverage, and the localization assumptions in their Assumption~2, it yields local $R$-linear convergence of appropriately selected proximal sequences. Their Theorem~1 also represents truncated operator values and inverse values as restricted sublevel sets of weakly convex functions; the inverse value at zero describes a localized zero set. Their Proposition~4 describes the resolvent of the truncated graph on its domain.

The compensation between a violation of firm nonexpansiveness and residual regularity has an established quantitative form. Lauster and Luke~\cite[Theorem~16, Lemma~17, and Corollary~19]{LausterLuke2021} use general regularity functions and summable comparison iterates to analyze fixed-point convergence. In the monotone proximal setting, Li and Mordukhovich~\cite[Theorem~7.3]{LiMordukhovich2012} derive convergence rates from H\"older metric subregularity. Thus, both nonlinear regularity and the use of regularity to compensate for an expansion are part of the existing theory. Our question is the exact relation between an input-anchored defect and an output residual bound for the same ordinary proximal step.

Alcantara's preprint~\cite{Alcantara2026}, first submitted on 16 November 2026, provides a close comparison. Proposition~3.4 and Remark~3.5 give step bounds and a local linear convergence estimate that overlap with DD1 below. The full paper treats degenerate preconditioned proximal methods, with reduced zero set $C^*(\zer A)$ and additional assumptions for identifying projected limits and recovering the original variables. Our comparison uses the scalar residual and anchored bounds in the ordinary proximal specialization; the graph and existence assumptions are treated separately.

Brito, Lara, and Yuan~\cite{BritoLaraYuan2026}, first submitted on 18 November 2026, address proximal existence directly. Under star monotonicity at a zero and alternative topological assumptions, their Theorem~3.1 proves global surjectivity of $I+cF$. Their Corollary~3.12 gives a surjectivity range for a quadratic radial violation. The proof first bounds homotopy solutions and then applies selection and fixed-point theorems. DD3 follows this argument in finite dimensions, imposing a bound on the entire operator graph relative to a nearest point in a specified closed subset of the zero set.

Throughout the paper, $H=\R^n$ has its Euclidean inner product and norm, $F:H\rightrightarrows H$, and $S$ is a nonempty closed subset of
\[
\zer F:=\{x\in H:0\in F(x)\}.
\]
We write
\begin{equation}\label{eq:intro-notation}
 d(x):=\dist(x,S),\qquad
 P_S(x):=\{p\in S:\norm{x-p}=d(x)\},\qquad
 r_F(y):=\inf_{v\in F(y)}\norm{v}.
\end{equation}
The infimum over an empty set is $+\infty$. For $\lambda>0$, the ordinary resolvent is
\[
 J_\lambda:=(I+\lambda F)^{-1},\qquad
 y\in J_\lambda(x)\ \Longleftrightarrow\ x-y\in\lambda F(y).
\]
The residual $r_F$ always uses the full value $F(y)$, even when an iteration uses a specified subset of the resolvent values.

To explain the two directions of DD1, consider an allowed proximal step $y\in J_\lambda(x)$ and suppose that
\begin{equation}\label{eq:intro-two-assumptions}
 d(y)\le\ell r_F(y),\qquad
 \ip{x-y}{y-p}\ge\Phi(d(x))
 \quad\text{for some }p\in P_S(x),
\end{equation}
where $\ell>0$. The first inequality measures regularity at the output. The second bounds the possible negative inner product relative to a nearest point to the input. On a distance interval $0<t\le R$, define
\begin{equation}\label{eq:intro-indices}
 L:=\frac{\ell}{\lambda},\qquad
 \sigma(t):=\frac{[-\Phi(t)]_+}{t^2},\qquad
 \sigma_R:=\sup_{0<t\le R}\sigma(t),\qquad
 [a]_+:=\max\{a,0\}.
\end{equation}
We call $L$ the \emph{regularity index} and $\sigma(t)$ the \emph{hypomonotonicity index} at scale $t$. These numerical indices measure the residual bound and the permitted negative inner product in~\eqref{eq:intro-two-assumptions}; the latter concerns a nearest input anchor rather than all pairs of graph points. Smaller values impose stronger bounds. Both depend on the chosen proximal parameter, and operators with the same indices satisfy the same compensation test.

For a fixed regularity index, the first direction computes the admissible defect. DD1 gives uniform strict distance decrease under
\begin{equation}\label{eq:intro-threshold}
 \sigma_R<T(L),\qquad
 T(L):=
 \begin{cases}
 (1-L)/L^2,&0<L\le\tfrac12,\\[2pt]
 1/(2L^2),&L\ge\tfrac12.
 \end{cases}
\end{equation}
Conversely, for a prescribed finite defect bound $\sigma_R=\sigma$, the required regularity is expressed by
\begin{equation}\label{eq:intro-inverse}
 L<\Lambda(\sigma),\qquad
 \Lambda(\sigma):=
 \begin{cases}
 +\infty,&\sigma=0,\\[2pt]
 (2\sigma)^{-1/2},&0<\sigma\le2,\\[2pt]
 2/(1+\sqrt{1+4\sigma}),&\sigma\ge2.
 \end{cases}
\end{equation}
The formulas agree at their shared endpoints. DD1 also gives the corresponding inverse calculation for any prescribed contraction factor. The piecewise relation comes from retaining both the step bound and the squared distance inequality. Together they give a geometric distance estimate and a step bound proportional to input distance. Under proximal coverage, absorption at zero distance, and the stated localization assumptions, the steps are summable and the iterates converge to a point of $S$.

The comparison with existing estimates is correspondingly specific. Alcantara~\cite[Remark~3.5]{Alcantara2026} explicitly obtains the condition $2\sigma L^2<1$ in the ordinary proximal specialization, and the step estimate in the proof of Proposition~3.4 supplies another ingredient of~\eqref{eq:intro-threshold}. Retaining both estimates adds the parameter region
\[
 0<L<\tfrac12,\qquad
 \frac{1}{2L^2}\le\sigma<\frac{1-L}{L^2}.
\]
The sharpness statement below allows the operator and the closed target set to vary; an individual operator may admit stronger estimates.

DD2 keeps the same proximal relation and replaces the linear residual bound by $d(y)\le\psi(r_F(y))$. For a general signed function $\Phi$, it computes a piecewise condition ensuring the prescribed decrease $d(y)\le\varphi(d(x))$. Examples include sublinear trajectories of finite length and a trajectory converging to a point with infinite length. We then obtain finite length from comparison series or an integrable potential. DD3 establishes global proximal existence and supplies the coverage used in these local convergence results. The final section starts from data on the original $F$ and determines a common admissible range from which parameters may be chosen at each iteration.

Section~\ref{sec:dd1} gives the two DD1 calculations and their geometric proof. Section~\ref{sec:dd2} develops compensation for general functions and the associated length estimates. Section~\ref{sec:dd3} proves the following existence lemma and combines it with DD1 and DD2. Section~\ref{sec:parameters} treats proximal parameter selection.

\begin{lemma}[Global proximal existence]\label{lem:existence}
Let $F:\R^n\rightrightarrows\R^n$ have nonempty values everywhere, let $S\subseteq\zer F$ be nonempty and closed, and let $c>0$. For each $\mu\in(0,c]$, let $\Phi_\mu:[0,\infty)\to\R$. Assume that, for every $\mu\in(0,c]$, every $y\in\R^n$, and every $u\in F(y)$, setting
\[
 z=y+\mu u,\qquad t=d(z),
\]
there exists $p_z\in P_S(z)$ such that
\begin{equation}\label{eq:global-anchor}
 \ip{z-y}{y-p_z}\ge\Phi_\mu(t).
\end{equation}
Assume also that
\begin{equation}\label{eq:global-defect}
 D_c(R):=\sup_{\substack{0<\mu\le c\\0\le t\le R}}
 [-\Phi_\mu(t)]_+<\infty
 \qquad\text{for every finite }R\ge0.
\end{equation}
Suppose that either
\begin{enumerate}
 \item $F$ is lower semicontinuous with nonempty closed convex values; or
 \item $F$ is upper semicontinuous with nonempty compact convex values.
\end{enumerate}
Then $\ran(I+cF)=\R^n$.
\end{lemma}
